%% ma: L10-B14-0053
%% lop: 10
%% chuong: 5
%% bai: 14
%% dang: 10.14.9
%% loai: TL
%% muc: VD
%% dap_an: "-"
%% nguon: "chuyen-de-cac-so-dac-trung-cua-mau-so-lieu-khong-ghep-nhom-toan-10-chuong-trinh-moi.pdf (GV. Phan Nhật Linh) (P05), Bài 03 Dạng 2 Bài tập tự luận 13"
%% kien_thuc: "Lớp 10 Bài 13 (số trung bình), Bài 14 (phương sai, độ lệch chuẩn)"
%% trang_thai: chinh-thuc
%% ghi_chu: "Tự duyệt theo yêu cầu giáo viên 22/10/2026."
\begin{ex}
Hai lớp 10A và 10B cùng làm bài thi môn Ngữ văn theo cùng một đề. Kết quả được trình bày ở hai bảng tần số sau:
\begin{center}
\begin{tabular}{|c|c|c|c|c|c|c|}\hline
Điểm thi 10A & $5$ & $6$ & $7$ & $8$ & $9$ & $10$\\\hline
Tần số & $6$ & $2$ & $10$ & $5$ & $6$ & $1$\\\hline
\end{tabular}
\end{center}
\begin{center}
\begin{tabular}{|c|c|c|c|c|}\hline
Điểm thi 10B & $6$ & $7$ & $8$ & $9$\\\hline
Tần số & $9$ & $9$ & $9$ & $3$\\\hline
\end{tabular}
\end{center}
Lớp nào có kết quả thi môn Ngữ văn đồng đều hơn?
\loigiai{$\overline x_A=\dfrac{216}{30}=7{,}2$; $\overline x_B=\dfrac{216}{30}=7{,}2$.

$s_A^2=\dfrac{6(-2{,}2)^2+2(-1{,}2)^2+10(-0{,}2)^2+5(0{,}8)^2+6(1{,}8)^2+(2{,}8)^2}{30}=\dfrac{62{,}8}{30}\approx2{,}09$.

$s_B^2=\dfrac{9(-1{,}2)^2+9(-0{,}2)^2+9(0{,}8)^2+3(1{,}8)^2}{30}=\dfrac{28{,}8}{30}=0{,}96$.

Vì $s_B^2<s_A^2$ nên lớp 10B có kết quả đồng đều hơn.}
\end{ex}
